\documentclass[11pt,a4paper]{article}
\usepackage[margin=24mm]{geometry}
\usepackage{fontspec}
\setmainfont{DejaVu Serif}[ItalicFont={DejaVu Serif},ItalicFeatures={FakeSlant=0.18},BoldItalicFont={DejaVu Serif Bold},BoldItalicFeatures={FakeSlant=0.18}]
\setmonofont{DejaVu Sans Mono}[Scale=0.8]
\usepackage{babel}
\babelprovide[main,import]{ukrainian}
\usepackage{amsmath,amssymb,booktabs,longtable,graphicx,microtype}
\usepackage[unicode,hidelinks]{hyperref}
\setlength{\emergencystretch}{3em}
\title{Мережа міст і громад у масштабі країни\\для двомасштабного методу}
\author{Робоче пояснення до основної SMPT-статті}
\date{16 жовтня 2026 року}
\begin{document}
\maketitle
\noindent Тут будуємо окремий великий приклад: мережу комун Франції на основі офіційних потоків між місцем проживання та роботи. Комуна --- адміністративна громада, яка може бути великим містом або малим населеним пунктом. Документ пояснює вибір даних, їх перетворення на граф і обмеження такого входу методу; це не повна журнальна стаття.

\section{Які два рівні обираємо}
Вузлом повної мережі є комуна, а спільнотою --- департамент. Тоді локальний граф описує зв'язки між комунами всередині департаменту; макрограф зберігає сумарні ваги між департаментами. Це дозволяє використати географічну мережу на десятки тисяч вузлів без вигадування контактів окремих людей.

Зміст стану тут змінюється порівняно з людиноцентричною моделлю: $I$ означає, що комуна вже охоплена процесом, а не що всі її жителі інфіковані. Такий SI-процес є моделлю просторового поширення між територіями. За цими даними не можна відновити кількість інфікованих людей у кожному місті. Пізніше окремо пояснимо, що потрібно для моделі, де локальними вузлами є люди, а спільнотами --- міста.

\section{Первинні джерела і географія}
Беремо \href{https://www.insee.fr/fr/statistiques/8582949}{INSEE, Mobilités professionnelles en 2022}, повний CSV потоків. Електронна таблиця з тієї самої сторінки містить лише потоки понад 100, тому вона непридатна як повна мережа. Оцінки стосуються 2022 року, але опубліковані в географії на 1 січня 2025 року. Для узгодження кодів використовуємо саме \href{https://www.insee.fr/fr/information/8377162}{Code officiel géographique 2025}.

Область даних --- Франція без Майотти. Париж, Ліон і Марсель подані через муніципальні округи, які треба звести до батьківських комун. До обробки не припускаємо, що всі коди місця роботи є французькими комунами: закордонні та інші коди перевірятимемо за офіційним довідником.

Джерело оцінює кількість зайнятих людей віком від 15 років за парою місць проживання та роботи; це не лічильник поїздок і не вимірювання фізичних контактів. INSEE застерігає, що потоки, менші за 200, слід розглядати як порядки величини. Зберігатимемо всі додатні потоки, не перетворюючи цю застережну межу на довільний поріг вилучення ребер.

Атрибуція: Source: Insee. Випуск потоків --- 26.06.2025; географічний довідник --- на 01.01.2025. \href{https://www.insee.fr/fr/information/2008466}{Умови INSEE} передбачають Licence Ouverte / Open Licence 2.0, якщо окремо не зазначено інше. URL і SHA-256 оригіналів зберігаються в \path{sources.json}.

\section{Що знайшла перевірка сирих даних}
Повний файл містить 986\,967 рядків. Первинний ключ --- пара кодів місця проживання та роботи; його дублікатів немає. У потрібних полях немає пропусків, усі оцінки потоку додатні та скінченні. Значення дробові, бо це статистичні оцінки, а не перелік ідентифікованих людей. Не округлюємо їх до цілих і не створюємо по одному ребру на кожного працівника.

Довідник містить 34\,875 комун, із яких 17 належать Майотті. Після узгодженого обмеження області залишаються 34\,858 комун у 100 департаментах. Додаємо всі ці комуни, зокрема ті, для яких немає міжкомунальних потоків. Відсутність спостереженого зв'язку не доводить відсутності всіх реальних контактів.

Для місць роботи є 29\,496 рядків поза вибраною областю. Їхня сумарна оцінка --- 497\,323.275 працівника. З них 310 рядків ведуть до Майотти; 29\,186 містять коди, яких немає в таблиці комун і муніципальних округів COG. Серед останніх є закордонні та невизначені місця роботи. Їх не перейменовуємо на французькі комуни й не додаємо як однакові міські вузли. Усі ці рядки збережено окремо в \path{outside_scope_flows.csv.gz} з причиною відокремлення. Відсікання вилучає спостережені зв'язки мешканців вибраної області з місцями роботи поза нею. Напрямлених потоків із зовнішніх місць проживання цей файл не містить, тому зовнішнє занесення та взаємодія із зовнішнім резервуаром потребували б додаткових даних і моделі.

Для французьких місць проживання всі коди узгодилися з довідником. Муніципальні округи Парижа, Ліона та Марселя зводимо до їхніх батьківських комун за полем \texttt{COMPARENT}, без ручного вгадування коду міста. Після такого зведення кілька різних первинних пар можуть стати однією парою комун. Це законна агрегація, а не дублікат джерела: її ваги підсумовуються.

У вибраній області 957\,471 первинний рядок перетворюється на 919\,779 напрямлених ненульових пар, включно з діагоналлю. Облік сум збережено:
\begin{align*}
28\,382\,870.874 &= 27\,885\,547.599 + 497\,323.275,\\
27\,885\,547.599 &= 9\,647\,898.373 + 18\,237\,649.225.
\end{align*}
У першому рядку повна оцінка розділена на внутрішню область та відокремлені потоки. У другому --- внутрішня область на роботу в тій самій комуні та роботу в іншій комуні. Показані числа округлено для читання; перевірки виконуються за збереженими значеннями подвійної точності.

\paragraph{Чому не відкидаємо малі потоки.}
969\,071 рядок має оцінку меншу за 200. Це переважна частина первинних пар. Попередження про їхню точність не означає, що їх немає. Порогове видалення могло б розірвати сполучення й змінити маршрути, а разом із ними --- похибку двомасштабного наближення. Основний граф зберігає всі додатні потоки; пороговий граф був би окремою версією з іншим змістом.

\section{Від напрямленого потоку до графа для наявного коду}
Нехай $F_{uv}$ --- оцінка кількості зайнятих людей, які проживають у комуні $u$ і працюють у комуні $v$, після географічного зведення. Матриця $F$ напрямлена та зберігається окремо, разом із діагоналлю. Наявні локальний і повний симулятори працюють із простим неорієнтованим графом. Для узгодженого входу визначимо
\[
a_{uv}=\frac{F_{uv}+F_{vu}}{s},\quad u\ne v,
\qquad a_{uu}=0,\qquad s=1000.
\]
Це симетричний показник зв'язності двох комун на основі зайнятості. Він \emph{не є виміряною кількістю контактів} і не означає, що потоки в обидва напрямки однакові. Симетризація є явним модельним рішенням; первинний напрямок не втрачений у файлі $F$, але не використовується базовим графом $A$.

Наприклад, якщо $F_{uv}=1200$, а $F_{vu}=300$, обидва напрямки поширення на проксі-графі отримують вагу $a_{uv}=a_{vu}=1.5$. Це не два фактично спостережені потоки по 1500 людей. Число 1000 лише задає масштаб ваг. Якби використовували $s=1$, для тотожних інтенсивностей потрібно було б зменшити коефіцієнт $\beta$ у 1000 разів.

Для обраного просторового SI-процесу сприйнятлива комуна $v$ отримує подію охоплення з інтенсивністю
\[
q_v(x)=\beta\sum_{u:x_u=I}a_{uv}.
\]
Якщо час виражений у днях, $\beta$ є інтенсивністю на одиницю проксі-ваги за день. Дані не задають її значення. Не можна назвати $F$ кількістю щоденних поїздок, поділити її на 365 чи без додаткової моделі прирівняти $\beta$ до ймовірності передачі під час однієї поїздки.

\paragraph{Чому прибираємо діагональ саме тут.}
Після зведення округів $F_{uu}$ містить також переміщення між округами того самого міста. Вони мають значення для людиноцентричної моделі всередині міста. Але в поточній просторовій мережі місто --- один бінарний вузол: самоконтакт не може перевести його із $S$ в $I$, якщо воно вже є джерелом $I$. Тому ці оцінки зберігаються як атрибут і в $F$, а не як петлі $A$. Це зміна роздільності, яка не відновлює внутрішнє зараження жителів.

\section{Як отримати локальні графи та матрицю методу}
Нехай $V_i$ --- комуни департаменту $i$. Визначаємо
\[
G_i=G[V_i],\qquad N_i=|V_i|,\qquad
W_{ij}=\sum_{u\in V_i,\,v\in V_j}a_{uv}\quad(i\ne j),\qquad W_{ii}=0.
\]
Повний граф використовує $A$, а двомасштабна модель --- ті самі локальні ребра й $W$, побудовану з $A$. Макроінтенсивність базового методу тоді
\[
\lambda_{ij}=\beta W_{ij}\frac{I_i}{N_i}\frac{S_j}{N_j}.
\]
За події $i\to j$ одержувач обирається рівномірно серед сприйнятливих комун $V_j$:
\[
\Pr(v\mid i\to j,x)=\frac{1}{S_j},\qquad v\in V_j,\quad x_v=S.
\]
Конкретний вузол-джерело базовий макрошар не обирає й не передає: він використовує лише $I_i/N_i$. Ця втрата розташування джерел і правило одержувача є частинами наближення разом з формулою інтенсивності.

Число $I_i$ тут рахує охоплені комуни, а $N_i$ --- усі комуни департаменту. Це не частка інфікованих працівників або мешканців. Велика й мала комуна додають по одному до цього лічильника; вплив відмінності розмірів частково присутній у вагах, але це не замінює модель населення всередині міста. Абсолютні потоки можуть посилювати тиск на великі за зайнятістю комуни та від них, тоді як кожна охоплена комуна рахується однією подією. Результати такої моделі відображатимуть і розмір зайнятої сукупності, і мережеве положення; це не нормований на населення територіальний ризик.

Матриця напрямлених потоків департаментів теж збережена, проте лише як описова таблиця. Не слід подавати її замість $W$ у порівняння з симетричним повним графом: тоді моделі описували б різні інтенсивності ще до двомасштабного наближення. Також не можна залишати в $W$ внутрішньодепартаментні суми: вони вже враховані локальними графами й не є макроімпортом.

Для збереження кодів і сум виконані перевірки:
\[
\sum_iN_i=34\,858,\qquad
|E|=|E_{\rm intra}|+|E_{\rm inter}|,
\]
\[
\frac12\sum_{i,j}W_{ij}=\sum_{\{u,v\}\in E_{\rm inter}}a_{uv},\qquad
s\sum_{u<v}a_{uv}=\sum_{u\ne v}F_{uv}.
\]
Остання рівність означає, що симетризація не подвоює загальну вагу, коли кожне неорієнтоване ребро рахується один раз. Окремо перевірено симетрію $A$ та $W$, нульові діагоналі і точну відповідність рядків $W$ списку департаментів.

\section{Яку мережу фактично отримано}
\begin{table}[htbp]
\centering
\caption{Основний підготовлений граф. Підрахунки є властивостями збережених даних, а не результатами епідемічних запусків.}
\begin{tabular}{lr}
\toprule
Показник & Значення\\
\midrule
Комуни (вузли повного графа) & 34\,858\\
Неорієнтовані зважені ребра & 743\,711\\
Департаменти (спільноти) & 100\\
Ребра всередині департаментів & 448\,147\\
Ребра між департаментами & 295\,564\\
Частка ваги всередині департаментів & 73.18\%\\
Ненульові неорієнтовані пари департаментів & 4611 із 4950\\
Компоненти зв'язності повного графа & 168\\
Вузли найбільшої компоненти & 34\,691\\
Ізольовані комуни & 167\\
Департаменти з незв'язним локальним графом & 56\\
Розміри спільнот: мінімум / медіана / максимум & 1 / 322.5 / 887\\
\bottomrule
\end{tabular}
\end{table}

Велика кількість ребер не означає майже повний граф: медіанний ступінь комуни дорівнює 20. Найбільший ступінь --- 8856. Мережа неоднорідна і за кількістю сусідів, і за вагами. Приклади найсильніших зв'язків наведено нижче. Це оцінки пов'язаності через зайнятість, а не щоденні транспортні рейси.

\begin{table}[htbp]
\centering\small
\caption{Вісім найбільших симетричних міжкомунальних зв'язків. Значення до ділення на $s=1000$; таблиця автоматично сформована зі збереженого графа.}
\input{figures/top_links.tex}
\end{table}

\begin{figure}[htbp]
\centering
\includegraphics[width=0.86\linewidth]{figures/department_links.pdf}
\caption{Підматриця 15 департаментів із найбільшою сумою зовнішніх ваг $\sum_jW_{ij}$. Колір показує $\log_{10}(1+1000W_{ij})$; він підкреслює як сильні, так і слабші зв'язки. Це лише візуальний піднабір, усі 100 департаментів збережені в $W$. Порядок визначено зовнішньою вагою, а не географічною відстанню; рисунок не є картою.}
\end{figure}

\begin{figure}[htbp]
\centering
\includegraphics[width=\linewidth]{figures/national_structure.pdf}
\caption{Ліворуч --- емпірична функція розподілу кількості сусідів серед 34\,691 неізольованої комуни; 167 ізольованих виключені лише з цієї логарифмічної візуалізації та залишені в мережі. Праворуч --- 12 департаментів із найбільшою кількістю комун. Розподіл описує вузли одного графа; він не є розподілом часів зараження.}
\end{figure}

\subsection*{Асиметрія напрямлених потоків}
Для напрямлених міжкомунальних потоків обчислено
\[
\eta=\frac{\sum_{u\ne v}|F_{uv}-F_{vu}|}{2\sum_{u\ne v}F_{uv}}
=0.6224.
\]
Нуль означав би повну симетрію. Значення показує значну напрямлену неоднорідність. Воно не є похибкою $A$ чи результатом порівняння симуляторів: симетричний проксі свідомо описує інший механізм, ніж односторонній рух. Саме тому первинна $F$ збережена окремо для можливого подальшого розвитку моделі.

Уточнення щодо порога 200: нижче нього лежать 98.19\% первинних рядків і 48.63\% повної ваги оцінених працівників. Просте відкидання цих рядків прибрало б майже половину ваги. Потоки поза вибраною географічною областю становлять 1.75\% повної ваги; це інше обмеження, яке не слід змішувати з порогуванням.

\section{Що така структура означає для похибки методу}
Мостові вузли існують у 30\,899 комун, але зовнішні ваги розподілені нерівномірно. Визначимо для комуни $u$ її зовнішню силу
\[
b_u=\sum_{v:\,d(v)\ne d(u)}a_{uv}.
\]
349 комун із найбільшими $b_u$, тобто верхній 1\% усіх вузлів із округленням угору, несуть 60.57\% сумарної зовнішньої сили. Це показник концентрації на кінцях ребер. Він не означає, що ці комуни завжди інфікуються першими.

Агрегована модель замінює реальний розподіл цієї ваги рівномірним доступом усіх вузлів спільноти. Якщо зараження ще не досягло комун із великою зовнішньою силою, макроінтенсивність може суттєво відрізнятися від точної. Якщо воно досягло їх рано, напрямок відмінності може бути іншим. Структура дає механізм для перевірки; вона не визначає наперед середнє зміщення часів приходу.

\paragraph{Ізольовані комуни.}
Повний граф складається з однієї великої компоненти та 167 ізольованих вузлів. Якщо ізольована комуна початково сприйнятлива, повний SI-процес на цьому графі не зможе її охопити. Рівномірний макроімпорт у її департамент може обрати її одержувачем, хоча фізичного ребра проксі-графа до неї немає. Це можливе порушення маршрутів самого наближення, а не помилка числового інтегрування.

Тому час охоплення всіх 34\,858 комун не є придатною скінченною метрикою повного процесу за звичайного одиничного початкового зараження. Для фіксованого початкового вузла $v_0$ визначаємо контрольну множину $C(v_0)$ як його компоненту зв'язності в $A$. Розподіли приходу порівнюємо на тих самих цільових комунах із $C(v_0)$ для обох методів. Комуни поза $C(v_0)$ позначаємо як структурно недосяжні для повного процесу; це інша категорія, ніж права цензура в час $T_{\rm end}$ для ще не охопленої досяжної комуни. Охоплення поза $C(v_0)$ у наближенні слід звітувати окремо як порушення досяжності, а не трактувати як звичайний завершений прихід. Це правило метрик, а не видалення вузлів із динаміки моделі. Для кількох початкових вузлів беруть об'єднання їхніх контрольних компонент.

\paragraph{Незв'язні локальні графи.}
У 56 департаментів локальна мережа незв'язна. Одна внутрішня інфекція не обов'язково охопить весь департамент без нових імпортів. Не можна автоматично використати часовий розкид зараження всіх його комун, як для зв'язного синтетичного графа: частина локальних часів може бути нескінченною. Тому безумовний $\sigma_T$ для повного локального охоплення може не існувати як скінченна величина.

Також показник $\rho$ головної статті не можна механічно обчислити як добуток нашого коефіцієнта $\beta$ та локального часового розкиду. Якщо замінити $s$ на $cs$, ваги $A$ і $W$ зменшаться в $c$ разів; заміна $\beta$ на $c\beta$ залишить усі інтенсивності та розподіли часу незмінними, але змінить $\beta\sigma_T$ у $c$ разів. Отже, на цій зваженій мережі потрібен інваріантний до одиниць ваг ефективний масштаб зовнішньої інтенсивності та визначення локальних часів з урахуванням незв'язності. Такого показника тут не визначали й не оцінювали. Поріг $\rho<0.25$ із синтетичних тестів головної статті не є гарантією для цієї мережі.

\paragraph{Що буде чесним порівнянням.}
Full та двомасштабний метод мають отримати той самий проксі-граф через узгоджені $A,G_i,W$, однакову сукупність вузлів і початкову умову, спільні горизонт та одиницю часу. Для кожного департаменту можна оцінювати час першого охоплення, ймовірність приходу до горизонту і частку охоплених комун. Окремо потрібні ознаки допустимості переходів та перевірка кроку синхронізації.

Не можна трактувати різницю між напрямленим транспортним процесом на $F$ та симетричним процесом на $A$ як похибку двомасштабного наближення. Тут є принаймні три різні рішення: вибір просторових вузлів, заміна зайнятості проксі-зв'язками та агрегація між департаментами. Синхронізація додає ще окрему числову похибку. У поточній підготовці часи зараження й ці похибки не оцінювалися.

\section{Як зберігати країну без зайвих щільних матриць}
Матриця на $34\,858\times34\,858$ елементів типу float64 потребувала б 9.05~GiB лише для одного масиву. У розрідженому CSR-поданні масиви $A$ займають 17.16~MiB; це обсяг її значень та індексів у пам'яті, а не оцінка повної пам'яті симулятора. NetworkX і журнали подій можуть потребувати значно більше.

Щільною зберігаємо лише $W$ розміру $100\times100$: 80\,000 байтів. Цей вибір відповідає двом рівням моделі. Не слід подавати величезну матрицю комун як макровхід поточного рушія: він має кілька щільних робочих масивів і обробляє всі пари спільнот. Також завантажувач передає нейтральні \texttt{alphas=1}, тому немає потреби обчислювати щільні спектральні розклади локальних графів для базової моделі.

Усі канонічні результати лежать у \path{prepared/}:
\begin{itemize}
\item \path{nodes.csv}: цілочисловий ID, код і назва комуни, департамент, регіон, ознаки ступеня та границі. Кількість зайнятих із рядків джерела не називається населенням комуни.
\item \path{departments.csv}: точний порядок спільнот, їхні розміри, компоненти зв'язності та ваги.
\item \path{F_domestic.npz}: напрямлені оцінки, включно з роботою в тій самій комуні; \path{A_proxy.npz}: симетричний граф без петель.
\item \path{W.npy}, \path{W.csv}: матриця методу з нульовою діагоналлю. Описові напрямлені потоки департаментів мають окремі файли з позначкою \texttt{descriptive}.
\item \path{geography_crosswalk.csv}: відображення кодів комун і округів на батьківську комуну.
\item \path{outside_scope_flows.csv.gz}: усі потоки поза областю, без прихованого видалення їх зі звітності.
\item \path{summary.json}, \path{components.csv}, \path{top_links.csv}: перевірки та числові результати; \path{manifest.json}: версії бібліотек, конфігурація й хеші всіх артефактів.
\end{itemize}
Незмінені оригінали лежать у \path{raw/}; \path{sources.json} зберігає URL, час завантаження, умови й хеші. \path{preparation_runs.jsonl} фіксує початок, завершення, помилки та пропуски готових результатів. Це журнал детермінованої підготовки, а не епідемічних експериментів.

\section{Команди відтворення та підключення до коду}
З кореня репозиторію:
\begin{verbatim}
.venv/bin/python paper/national_network_preparation/fetch_sources.py
.venv/bin/python paper/national_network_preparation/prepare_network.py
.venv/bin/python paper/national_network_preparation/analyze_network.py
\end{verbatim}
Завантажувач перевіряє наявні оригінали й пропускає повторне скачування. Підготовка та рисунки теж перевіряють хеші. Зміна коду, вхідних даних або масштабу ваг потребує нового каталогу результатів; старі артефакти не перезаписуються. Команди тестів та PDF-збірки наведені в README.

Для роботи з канонічними розрідженими даними:
\begin{verbatim}
from paper.national_network_preparation.prepare_network import load_sparse
nodes, departments, A, W = load_sparse()
\end{verbatim}
Для явного перетворення у формат наявних симуляторів:
\begin{verbatim}
from paper.national_network_preparation.prepare_network import load_networkx
micro, full, W, comm_nodes = load_networkx()
\end{verbatim}
Перетворення створює один повний NetworkX-граф і локальні індуковані представлення без копій усіх ребер. Зворотне відображення цілочислових ID у географічні коди задане в \path{nodes.csv}; коди з провідними нулями або літерами не перетворюються на числа.

Повний API приймає \texttt{full\_graph=full} та \texttt{comm\_nodes=comm\_nodes}. Двомасштабний --- \texttt{micro\_graphs=micro}, \texttt{W=W} та \texttt{full\_graph=full} для метаданих. Для базової перевірки встановлюють \texttt{alpha\_mode="off"} і \texttt{correction\_mode="off"}; коефіцієнти локальної й зовнішньої передачі мають бути узгоджені з тим самим $A$. Для SI одужання дорівнює нулю. Значення $\beta$, горизонт і крок тут не призначалися.

Перевірено саме завантаження повного національного графа та конструктори обох API-обгорток. Це не запуск їхніх циклів. Під час майбутнього великого запуску слід також контролювати розмір журналів: прямий оркестратор підтримує \texttt{dense\_logs=False}, тоді як публічна обгортка наразі цього параметра не відкриває. При підготовці ядро симулятора не змінювали.

\section{Якщо місто має бути спільнотою людей}
Побудований варіант має рівні \emph{комуна всередині департаменту}. Інша постановка --- \emph{людина всередині міста}, з містами на макрорівні --- також змістовна, але потребує інших входів. Для неї кількість жителів задає $N_i$, потрібні локальні мережі домогосподарств, шкіл і робочих місць та модель того, як мобільність створює контакти між людьми.

Матриця зайнятості $F$ може бути емпіричною основою такого макрозв'язку, але не визначає ні адресатів зараження, ні контактні тривалості, ні біологічну інтенсивність. Синтетичні локальні мережі можна генерувати з фіксованими seed і явно задокументованими припущеннями; це вже буде гібридна синтетично-емпірична модель. Для повного контрольного графа довелося б також задати конкретні міжміські контакти, яких у цих даних немає.

Для всіх 34\,858 комун як макроспільнот поточний щільний макрорушій потребуватиме зміни на розріджене представлення або іншого рівня географічного поділу. Наявна підготовка вже зберігає $F$, географічні коди й відокремлені потоки для такого продовження, але не видає модель комун за реконструйовану мережу десятків мільйонів людей.


\end{document}
